% ============================================================================
% Taller 2 — Modelos AR, MA y ARMA (Modulo II)
% Series de Tiempo (20948) · Corte II · Universidad El Bosque · 2026-II
%
% PLANTILLA DE ENTREGA. Se entrega el PDF compilado, por Brightspace.
%
% Como usarla:
%   1. Rellena la ficha de la portada. Los TRES digitos de verificacion son lo
%      primero que hay que comprobar: un ARMA perfectamente ajustado sobre la
%      serie equivocada se ve identico a uno correcto.
%   2. Rellena la hoja de cifras de la pagina 2 ANTES de escribir prosa.
%   3. Responde cada literal donde dice \marcador{...}. Los marcadores en gris
%      que queden sin sustituir se ven en el PDF: esa es su unica funcion.
%   4. Compila con pdflatex (o subela a Overleaf tal cual: no usa ningun
%      paquete raro ni ningun archivo externo) y entrega el PDF.
%
% Nombre del archivo:  T2_Apellido_XXX.pdf   (XXX = tus tres ultimos digitos)
% ============================================================================
\documentclass[11pt,a4paper,oneside]{article}

\usepackage[utf8]{inputenc}
\usepackage[T1]{fontenc}
\usepackage[spanish,es-nodecimaldot]{babel}
\usepackage{lmodern}
\usepackage{microtype}
\usepackage[a4paper,top=2.3cm,bottom=2.3cm,left=2.4cm,right=2.4cm]{geometry}

\usepackage{graphicx}
\usepackage{booktabs}
\usepackage{array}
\usepackage{tabularx}
\usepackage{caption}
\usepackage{enumitem}
\usepackage{fancyhdr}
\setlength{\headheight}{14pt}
\usepackage{titlesec}
\usepackage{amsmath,amssymb}
\usepackage{xcolor}
\usepackage{listings}
\usepackage{float}
\usepackage{hyperref}

% --- Colores institucionales (los mismos del material de Series de Tiempo) ---
\definecolor{ubosque}{RGB}{1,40,32}
\definecolor{naranja}{RGB}{255,102,0}
\definecolor{grisclaro}{RGB}{240,240,240}
\definecolor{grismarcador}{RGB}{130,130,130}

\hypersetup{colorlinks=true, linkcolor=ubosque, urlcolor=blue!70!black,
            pdftitle={Taller 2 - Series de Tiempo 20948}}

\pagestyle{fancy}
\fancyhf{}
\fancyhead[L]{\small\textit{Taller 2 --- Series de Tiempo (20948)}}
\fancyhead[R]{\small\textit{\leftmark}}
\fancyfoot[C]{\thepage}
\fancyfoot[R]{\small Universidad El Bosque}
\renewcommand{\headrulewidth}{0.4pt}
\renewcommand{\footrulewidth}{0.3pt}

\setlength{\parskip}{0.5\baselineskip}
\setlength{\parindent}{0pt}
\emergencystretch=3em
\titleformat{\section}{\normalfont\large\bfseries\color{ubosque}}{\thesection.}{0.5em}{}
\titleformat{\subsection}{\normalfont\normalsize\bfseries}{}{0em}{}
\captionsetup{font=small,labelfont=bf,labelsep=period}

\lstset{basicstyle=\ttfamily\footnotesize, breaklines=true, frame=single,
        rulecolor=\color{grisclaro!60!black}, columns=flexible,
        showstringspaces=false, xleftmargin=0pt, framexleftmargin=3pt}

% --- Los tres comandos de la plantilla ---
% Un hueco por rellenar. Se ve en gris y entre corchetes: si queda uno sin
% sustituir, salta a la vista en el PDF entregado.
% Ojo con la forma del color: `\textcolor{}{}` (comando) y NO
% `{\color{} ...}` (declaracion). Dentro de una celda de parrafo —las
% columnas X de tabularx— la declaracion arranca el parrafo con un whatsit y
% hunde su primera linea: la etiqueta de la fila queda arriba y su casilla
% media linea mas abajo. Aislado en un caso minimo: `{\itshape[x]}` alinea,
% `{\color{gris}[x]}` no. (Lección heredada de la plantilla de Estadistica
% Espacial; no cambiar sin reproducir el caso minimo.)
\newcommand{\marcador}[1]{\textcolor{grismarcador}{\textit{[#1]}}}

\newcommand{\tarea}[4]{%
  \section[#1 · #2]{#1 · #2 \normalfont\small\color{grismarcador}(#3 \% del escrito · #4 palabras)}%
  \markboth{#1 · #2}{}}

% Un literal: la letra en negrita y el rotulo corto. El enunciado COMPLETO
% esta en el taller; aqui solo va el rotulo, para no gastar dos paginas en
% copiar lo que el estudiante ya tiene delante.
\newcommand{\literal}[2]{\vspace{0.35\baselineskip}\noindent\textbf{(#1)}~\textit{#2}\par}

\begin{document}

% ============================================================================
% PORTADA — la ficha de variante
% ============================================================================
\thispagestyle{empty}
\begin{center}

{\large\textbf{UNIVERSIDAD EL BOSQUE}}\\[0.15cm]
{\normalsize Facultad de Ciencias --- Matemáticas y Ciencia de Datos}\\[0.1cm]
{\normalsize Series de Tiempo (20948) --- 2026-II --- Corte II}\\[0.9cm]

\rule{\textwidth}{1.2pt}\\[0.25cm]
{\LARGE\textbf{Taller 2}}\\[0.15cm]
{\large\textit{Modelos AR, MA y ARMA:\\ identificación, estimación y diagnóstico}}\\[0.25cm]
\rule{\textwidth}{1.2pt}\\[0.9cm]

{\large\textbf{Trabajo individual}}\\[0.15cm]
{\normalsize Escrito 60\,\% \quad · \quad Defensa oral 40\,\%}\\[1.0cm]

\end{center}

\begin{tabular}{rl}
\textbf{Nombre completo:}          & \marcador{tu nombre} \\[0.25cm]
\textbf{Número de documento:}      & \marcador{completo} \\[0.25cm]
\textbf{Tus tres últimos dígitos:} & \marcador{000 a 999} \quad $\rightarrow$ \quad
                                     \textbf{variante} \marcador{la misma cifra} \\[0.25cm]
\textbf{Tu serie A (T1, T4):}      & número \marcador{NN} \quad dígito de verificación \marcador{suma} \\[0.25cm]
\textbf{Tu serie B (T2):}          & número \marcador{NN} \quad dígito de verificación \marcador{suma} \\[0.25cm]
\textbf{Tu serie C (T3):}          & número \marcador{NN} \quad dígito de verificación \marcador{suma} \\[0.25cm]
\textbf{Fecha de entrega:}         & \textbf{martes 6 de octubre de 2026, 14:00} \\
\end{tabular}

\vspace{0.9cm}

\noindent\fcolorbox{ubosque}{grisclaro}{%
\begin{minipage}{0.95\textwidth}
\small
\textbf{Antes de escribir nada, comprueba los tres dígitos de verificación.} Cada uno es la suma de
los valores de esa serie, redondeada. Si alguno no da exactamente esa cifra, no estás trabajando
con tu variante, y un ARMA perfectamente ajustado sobre la serie equivocada se ve idéntico a uno
correcto. Todo lo que construyas encima estaría mal sin que nada avise.
\end{minipage}}

\vspace{0.5cm}

\noindent\fcolorbox{grisclaro!60!black}{white}{%
\begin{minipage}{0.95\textwidth}
\small
\textbf{Cómo se entrega.} Un solo PDF por \textbf{Brightspace}, compilado con \LaTeX{} a partir de
esta plantilla. Nombre del archivo: \texttt{T2\_Apellido\_XXX.pdf}, donde \texttt{XXX} son tus tres
últimos dígitos. Máximo \textbf{12 páginas} sin contar el anexo. Las tablas y figuras no cuentan
para los límites de palabras de cada tarea.

\smallskip
\textbf{El enunciado no está aquí.} Vive en el taller publicado, con tu variante resuelta al
escribir tu documento. Esta plantilla solo recoge las respuestas.
\end{minipage}}

\vfill
{\footnotesize\color{grismarcador}\today}

\newpage

% ============================================================================
% HOJA DE CIFRAS
% ============================================================================
\section{Hoja de cifras}
\markboth{Hoja de cifras}{}

Todas las cifras que el taller pide, juntas. \textbf{Rellénala antes de escribir prosa}: es lo que
permite que la discusión de cada tarea hable de números que ya están decididos, y es lo primero que
se revisa al calificar. Usa punto decimal y cuatro decimales donde la cifra los tenga.

\vspace{0.3cm}
\renewcommand{\arraystretch}{1.25}
\begingroup\setlength{\parskip}{0pt}
\begin{tabularx}{\textwidth}{@{}l X r@{}}
\toprule
\textbf{Tarea} & \textbf{Cifra} & \textbf{Valor} \\
\midrule
T1 & Serie A: $n$ y suma (dígito de verificación)                              & \marcador{~} \\
T1 & Banda del correlograma, $\pm 1.96/\sqrt{n}$                              & \marcador{~} \\
T1 & Rezago en que «corta» la ACF / la PACF (o «no corta»)                     & \marcador{~} \\
T1 & Orden $(p,q)$ que sugiere la tabla ACF/PACF                               & \marcador{~} \\
T1 & Orden que elige el AICc, y su AICc                                        & \marcador{~} \\
T1 & Orden que elige el BIC, y su BIC                                          & \marcador{~} \\
T1 & Modelo elegido, y $p$ de Ljung–Box de sus residuales (12 rezagos)         & \marcador{~} \\
T1 & ¿Hay un modelo más simple a menos de 2 unidades de AICc? Cuál            & \marcador{~} \\
\midrule
T2 & Serie B: suma (dígito de verificación)                                    & \marcador{~} \\
T2 & ARMA(3,3): mayor error estándar / error estándar del MA(1)                & \marcador{~} \\
T2 & ARMA(3,3): módulo de la raíz AR y de la raíz MA más próximas entre sí     & \marcador{~} \\
T2 & \texttt{ar()}: orden elegido, y razón $a_2/a_1$ de sus dos primeros coeficientes & \marcador{~} \\
T2 & Rejilla $\{0..3\}^2$: modelo que elige el AICc y el que elige el BIC      & \marcador{~} \\
T2 & AICc del ARMA(3,3) menos AICc del mínimo de la rejilla                    & \marcador{~} \\
\midrule
T3 & Serie C: suma (dígito de verificación)                                    & \marcador{~} \\
T3 & AR(2) sobre los datos crudos: $p$ de Ljung–Box (12 rezagos)               & \marcador{~} \\
T3 & Elasticidad amplitud–nivel: pendiente e intervalo del 95\,\%              & \marcador{~} \\
T3 & $\lambda$ de Guerrero                                                     & \marcador{~} \\
T3 & AR(2) en la escala correcta: $\hat\phi_1$, $\hat\phi_2$ y $\hat\sigma^2$  & \marcador{~} \\
T3 & Pseudo-periodo antes / después de transformar                             & \marcador{~} \\
T3 & $\hat\phi_1$ por Yule–Walker, CSS y ML, y el error estándar de ML         & \marcador{~} \\
T3 & Distancia $|\hat\phi_1^{\mathrm{YW}} - \hat\phi_1^{\mathrm{ML}}|$ en errores estándar & \marcador{~} \\
\midrule
T4 & Réplicas en que la tabla ACF/PACF recupera tu $(p,q)$ (de 200)            & \marcador{~} \\
T4 & Réplicas en que el AICc recupera tu $(p,q)$                               & \marcador{~} \\
T4 & Réplicas en que el ganador por AICc pasa Ljung–Box sin ser tu orden       & \marcador{~} \\
\bottomrule
\end{tabularx}
\endgroup

\vspace{0.45cm}
\noindent\textbf{T1 · la rejilla.} Copia aquí las filas de la rejilla que compites entre sí: el ganador por
AICc, el ganador por BIC y el modelo más simple que quede a menos de 2 unidades del ganador, si lo hay.

\vspace{0.25cm}
\begingroup\setlength{\parskip}{0pt}
\begin{tabularx}{\textwidth}{@{}c c r r r r X@{}}
\toprule
\textbf{$p$} & \textbf{$q$} & \textbf{AICc} & \textbf{BIC} & \textbf{$\hat\sigma^2$} & \textbf{LB $p$} & \textbf{¿Por qué está en esta lista?} \\
\midrule
\marcador{~} & \marcador{~} & \marcador{~} & \marcador{~} & \marcador{~} & \marcador{~} & \marcador{~} \\
\marcador{~} & \marcador{~} & \marcador{~} & \marcador{~} & \marcador{~} & \marcador{~} & \marcador{~} \\
\marcador{~} & \marcador{~} & \marcador{~} & \marcador{~} & \marcador{~} & \marcador{~} & \marcador{~} \\
\bottomrule
\end{tabularx}
\endgroup

\vspace{0.2cm}
{\small\color{grismarcador} T5 no lleva cifras propias: sus datos vienen dados en el enunciado y son
los mismos para todo el curso.}

\newpage

% ============================================================================
% LAS CINCO TAREAS
% ============================================================================
\tarea{T1}{La firma que no se deja leer}{25}{700}
\literal{a}{Qué orden propone la tabla ACF/PACF para tu serie A, y con qué seguridad: con las barras que cruzan la banda y con $n$, no con la regla de memoria.}
\marcador{tu respuesta}

\literal{b}{Qué eligen el AICc y el BIC. Si no coinciden, cuál y por qué. Si hay un modelo más simple a menos de 2 unidades del ganador, qué haces con él.}
\marcador{tu respuesta}

\literal{c}{El diagnóstico del modelo elegido. Si pasa, qué prueba eso y qué \textbf{no} prueba.}
\marcador{tu respuesta}

\literal{d}{Un compañero tiene otra serie donde el mismo modelo pasa el diagnóstico. ¿Vienen las dos series del mismo proceso? Responde con lo que Ljung–Box mide.}
\marcador{tu respuesta}

\tarea{T2}{Grande por si acaso}{25}{800}
\literal{a}{Sin ajustar nada más: qué hay en la salida del ARMA(3,3) que dice que sobra. Las raíces de $\phi(B)$ y de $\theta(B)$, y los errores estándar.}
\marcador{tu respuesta, con las raíces calculadas}

\literal{b}{El patrón de los coeficientes del \texttt{ar()}, escrito como fórmula, y de qué proceso son esos pesos. Por qué \texttt{ar()} no pudo encontrarlo.}
\marcador{tu respuesta}

\literal{c}{La rejilla $\{0..3\}^2$: qué eligen AICc y BIC y a qué distancia queda el ARMA(3,3). Qué pasaría si comparas $\hat\sigma^2$ en vez de AICc, y por qué sería un error.}
\marcador{remite a la hoja de cifras, y comenta aquí lo que las cifras no dicen}

\literal{d}{El analista tiene razón en que pasa Ljung–Box. Por qué eso no lo hace válido.}
\marcador{tu respuesta}

\tarea{T3}{El diagnóstico que pasa}{20}{700}
\literal{a}{Qué otra cosa tiene que cumplir un residual además de pasar Ljung–Box, y qué gráfico lo muestra. Nómbralo antes de calcular.}
\marcador{tu respuesta}

\literal{b}{La elasticidad amplitud–nivel por bloques, con su intervalo. Qué transformación pide, y qué dice Guerrero.}
\marcador{tu respuesta}

\literal{c}{El reajuste en la escala correcta: qué cambió y qué no cambió.}
\marcador{tu respuesta}

\literal{d}{Yule–Walker, CSS y ML en la escala correcta: cuánto se separa el que se separa, en errores estándar, por qué justo ese, y cuál reportarías.}
\marcador{tu respuesta}

\tarea{T4}{Simular y reconocer}{15}{500}
\literal{a}{En qué fracción de las 200 réplicas la tabla ACF/PACF recupera tu $(p,q)$, y en cuál lo recupera el AICc.}
\marcador{remite a la hoja de cifras}

\literal{b}{En qué fracción el ganador por AICc pasa Ljung–Box \emph{sin} ser tu orden.}
\marcador{remite a la hoja de cifras}

\literal{c}{Con esas tres cifras delante, vuelve a T1(a)–(c): qué cambias, o por qué tu decisión ya lo tenía en cuenta.}
\marcador{tu respuesta}

\tarea{T5}{Auditar a la IA}{15}{600}
\literal{a}{Clasifica las tres afirmaciones como correcta, a medias o falsa, con una cifra o un contraejemplo detrás de cada veredicto.}
\marcador{tu respuesta}

\literal{b}{La falsa, reescrita para que quede correcta sin cambiar de tema.}
\marcador{tu reescritura}

\literal{c}{La de «a medias»: dónde se rompe exactamente, qué parte es cierta, cuál no, y qué cifra lo demuestra.}
\marcador{tu respuesta}

\literal{d}{Una frase: qué tipo de pregunta contesta bien un modelo sobre este material y cuál contesta mal, apoyada en lo que te pasó durante el taller.}
\marcador{una frase}

\newpage
% ============================================================================
% DECLARACIÓN DE USO DE IA
% ============================================================================
\section{Declaración de uso de inteligencia artificial}
\markboth{Declaración de uso de IA}{}

El uso de IA en este taller es \textbf{libre y no hay que pedir permiso}. Lo que se califica —el
10\,\% del escrito, dimensión de comunicación y honestidad— es esta declaración: qué consultaste y
\textbf{cómo lo verificaste}. Declarar no penaliza; no declarar, sí.

Escribe una fila por tarea. Si en alguna no consultaste nada, escríbelo: «no consulté» es una
declaración válida y dejarla en blanco no lo es.

\vspace{0.3cm}
\renewcommand{\arraystretch}{1.4}
\begingroup\setlength{\parskip}{0pt}
\begin{tabularx}{\textwidth}{@{}l X X@{}}
\toprule
\textbf{Tarea} & \textbf{Qué consulté y a qué modelo} & \textbf{Cómo verifiqué la respuesta} \\
\midrule
T1 & \marcador{~} & \marcador{~} \\
T2 & \marcador{~} & \marcador{~} \\
T3 & \marcador{~} & \marcador{~} \\
T4 & \marcador{~} & \marcador{~} \\
T5 & \marcador{~} & \marcador{~} \\
\bottomrule
\end{tabularx}
\endgroup

% ============================================================================
% DECISIONES PARA LA DEFENSA
% ============================================================================
\section{Decisiones que sostendré en la defensa}
\markboth{Decisiones para la defensa}{}

La defensa vale el 40\,\%, y hay una regla que le da dientes: \textbf{una decisión entregada por
escrito que no se pueda sostener en la defensa se recalifica}. Escribe entre tres y cinco
decisiones tuyas —no resultados: decisiones— y de qué tarea salen.

\begin{enumerate}[leftmargin=*,itemsep=4pt]
  \item \marcador{decisión} \hfill \textcolor{grismarcador}{\small (tarea \marcador{Tn})}
  \item \marcador{decisión} \hfill \textcolor{grismarcador}{\small (tarea \marcador{Tn})}
  \item \marcador{decisión} \hfill \textcolor{grismarcador}{\small (tarea \marcador{Tn})}
  \item \marcador{opcional} \hfill \textcolor{grismarcador}{\small (tarea \marcador{Tn})}
  \item \marcador{opcional} \hfill \textcolor{grismarcador}{\small (tarea \marcador{Tn})}
\end{enumerate}

\newpage
% ============================================================================
% ANEXO — no cuenta para el limite de paginas
% ============================================================================
\section*{Anexo. Código ejecutado y entorno}
\markboth{Anexo}{}
\addcontentsline{toc}{section}{Anexo}

Pega aquí el código que ejecutaste de verdad para T1, T2, T3 y T4 —el de T4 con la semilla del
enunciado— y la salida de \texttt{sessionInfo()}. No se califica su elegancia: sirve para saber sobre qué serie trabajaste si
alguna cifra no cuadra.

\begin{lstlisting}[language=R]
# Pega aqui tu codigo
\end{lstlisting}

\begin{lstlisting}
# Pega aqui la salida de sessionInfo()
\end{lstlisting}

\end{document}
